% Generated draft: requires mathematical and source review.
\documentclass[11pt]{article}
\usepackage[margin=25mm]{geometry}
\usepackage{fontspec}
\setmainfont{Latin Modern Roman}
\usepackage{amsmath,amssymb,mathtools,unicode-math}
\setmathfont{Latin Modern Math}
\usepackage{enumitem,needspace,xurl,hyperref}
\setlength{\parindent}{0pt}
\setlength{\parskip}{5pt}
\setlength{\emergencystretch}{3em}
\newcommand{\N}{\mathbb{N}}
\newcommand{\Z}{\mathbb{Z}}
\newcommand{\Q}{\mathbb{Q}}
\newcommand{\R}{\mathbb{R}}
\newcommand{\C}{\mathbb{C}}
\newcommand{\F}{\mathbb{F}}
\newcommand{\A}{\mathbb{A}}
\newcommand{\PP}{\mathbb{P}}
\newcommand{\E}{\mathbb{E}}
\newcommand{\eps}{\varepsilon}
\newcommand{\dd}{\,\mathrm{d}}
\begin{document}
\section*{UnsolvedMath ID: 30001980}
\textbf{Source problem code:} catalogue entry
\subsection{Definitions and mathematical statement}
Ballistic annihilation refers to particle dynamics in which particles move with constant velocities between collisions and annihilate upon collision. The supplied record does not specify the ambient space, particle geometry, collision criterion, or treatment of simultaneous collisions. Deterministic dynamics with random initial data means that, once an initial configuration is chosen, its evolution is determined: schematically,
\[
X(t)=\Phi_t(X_0),
\]
where \(X_0\) is a random initial configuration and \(\Phi_t\) is the evolution of a specified particle model. This notation describes the terminology only; neither the configuration space, the initial probability law, nor the evolution maps are supplied.

The retained source question asks whether continuum equations, such as the Boltzmann equation, can be rigorously derived as scaling limits of deterministic many-particle dynamics, particularly ballistic annihilation with random initial data. A precise convergence statement would require a specified microscopic model, initial distribution, scaling regime, macroscopic observable, limiting equation and initial condition, time interval, and mode of convergence. None of these is fixed by the supplied question, so no particular Boltzmann equation or convergence assertion is asserted here.

The supplied record explicitly marks the statement as \texttt{unrecoverable}. Its statement-verification note says that the official contribution uses the continuum-limit question as motivation and then derives its announced Boltzmann scaling result; the extraction is not a separately retained open problem, and no clean statement is promoted. Consequently, the retained question is an incomplete research direction rather than a well-defined conjecture. The supplied \texttt{open} label does not resolve this formulation issue or establish the status of any specific scaling-limit theorem.
\subsection{Short English statement}
The record asks about deriving continuum equations from deterministic particle dynamics with random initial data, especially ballistic annihilation. It supplies neither a complete model nor a precise limiting assertion. Its verification note also identifies the question as motivation for a reported scaling result, rather than a separately retained open problem. A precise theorem or conjecture cannot be recovered from this record alone.
\subsection{Sources}
\begin{itemize}
\item Supplied local UnsolvedMath registry record, permanent ID \(30001980\), source problem code \texttt{catalogue entry}, title ``Ballistic-Annihilation Scaling-Limit Report.'' Local registry paths: \texttt{data/unsolvedmath/problems.json} and \texttt{catalog/problems.json}. The supplied record retains the original question for provenance and has no clean statement.
\item Supplied report citation: \emph{Interplay of Analysis and Probability in Physics} (2012), Oberwolfach Reports, pp.~281--349, DOI \texttt{10.4171/owr/2012/06}. Source URL and supplied literature-source link, labelled ``original Oberwolfach report'': \url{https://doi.org/10.4171/owr/2012/06}. No individual proposer or contribution-level page reference is supplied.
\item Literature assessment dated exactly \texttt{2026-08-21}: the supplied curation labels the record open but describes it as a truncated setup for a Boltzmann-limit problem, not a complete claim that can be marked solved or disproved. This is supplied dated curation, not an independently verified literature assessment.
\item The supplied statement-verification note says that the official contribution presents the continuum-limit question as motivation and derives the announced Boltzmann scaling result; it does not retain a separate open problem. The record was updated at \texttt{2026-09-16T12:06:53Z}. These source-verification statements are supplied metadata and have not been independently checked.
\end{itemize}
\end{document}
